\section{Universal collections}\label{sec:universal}
Universal sets for four-colour cubes were studied in~\cite[Section~6]{BES}.
A collection is \emph{universal} if, for every five-colour target variety,
eight of its copies can form a corner solution of that variety. The chosen
copies and their placements may depend on the target. The palette is fixed,
and only proper rotations are allowed, as before.

\begin{theorem}\label{thm:universal}
The minimum size of a universal collection is $15$, attained by fifteen
distinct adjacent-repeat varieties. Requiring only the sixty adjacent-repeat
targets already forces this minimum.
\end{theorem}
\begin{proof}
For each colour $r$, universality requires all four bicoloured corner types
$b(r,p)$, $p\ne r$. Only adjacent-repeat sources repeating $r$ can supply
them, and each such source supplies exactly two. If there were at most two
such copies, covering all four types would force exactly two copies with
disjoint pairs of types. A target whose bicoloured pair equals either pair
would then require two copies from a pool of one. Thus at least three
adjacent-repeat sources are needed for each of the five repeated colours,
giving the lower bound $15$. Any universal collection of size fifteen has
exactly three adjacent-repeat sources of each repeated colour and no
opposite-repeat sources.

For the upper bound, use the directed-corner notation of
Section~\ref{sec:support} and read the colours modulo $5$. For every $r$ and
$p\notin\{r,r+1\}$, take the positive adjacent-repeat variety with repeated
colour $r$ and smaller pair $\{r+1,p\}$. These are fifteen distinct varieties.
For each $r$, their bicoloured pairs form a three-edge star on the other
four colours, so any two required bicoloured corners can be assigned to
distinct sources.

Represent each source by its smaller pair in an auxiliary multigraph on the
five colours. Pairs consecutive in the colour cycle occur once and the other
pairs twice. Every vertex therefore has degree six, and every pair has
multiplicity at most two. A source supplies the tricoloured corner $uv$
exactly when its pair contains $u$ but not $v$. Hence one arc has at least
four neighbouring sources. Two distinct arcs have at least six: if their
tails agree, their neighbourhood union contains all six incident sources;
otherwise each has at least four neighbours and their intersection has at
most two.

The tricoloured corners of any target have at most four tails and at most
three arcs per tail. A set of at least five such corners therefore contains
two arcs with the same tail and one with a different tail. Writing $u$ for
the common tail and $w\ne u$ for the third tail, a source common to the
six-source star and the third neighbourhood must have smaller pair $\{u,w\}$.
Their intersection therefore has size at most two, and their neighbourhood
union has size at least $6+4-2=8$. Thus a set of $k$ tricoloured target corners
has at least
\[
\begin{array}{c|ccc}
 k&1&2\text{--}4&5\text{--}8\\\hline
 \text{neighbouring sources}&4&6&8
\end{array}
\]
sources. For an opposite-repeat target these bounds imply every Hall
inequality. For an adjacent-repeat target, a Hall set with $k\ge1$
tricoloured corners has at most $k+2$ positions, and $k\le6$; the same bounds
again suffice. Sets containing only bicoloured corners were handled by the
star construction. Hall's theorem completes every target.
\end{proof}

Both bounds are independent of the certified fifteen-cube corner theorem.


\begin{samepage}
\begin{corollary}\label{cor:universal-structure}
Every minimum universal collection consists of fifteen distinct adjacent-repeat
varieties. For each repeated colour, their three smaller pairs form a spanning
tree on the other four colours.
\end{corollary}
\begin{proof}
There are exactly three adjacent-repeat sources of each repeated colour $r$.
Their smaller pairs must cover the other four colours. If two pairs coincide,
the third must be disjoint from them; a target requiring the two bicoloured
corners of that third pair would have only one supplier. Thus the three pairs
are distinct and cover four vertices, so they form a spanning tree. In
particular, the sources are distinct varieties. The construction above chooses
a star for each of these trees.
\end{proof}
\end{samepage}
